% Part: second-order-logic
% Chapter: syntax-and-semantics
% Section: semantic-notions

\documentclass[../../../include/open-logic-section]{subfiles}

\begin{document}

\olfileid{sol}{syn}{sem}
\olsection{مفاهیم معناشناختی}

\begin{explain}
مفاهیم منطقی محوریِ \emph{اعتبار}، \emph{نتیجهٔ منطقی} و
\emph{ارضاشدنی‌بودن} در منطق مرتبهٔ دوم همان‌گونه تعریف می‌شوند که در
منطق مرتبهٔ اول، جز اینکه رابطهٔ ارضای مبنا اکنون رابطهٔ مربوط به
!!{formula}s مرتبهٔ دوم است. البته یک !!{sentence} مرتبهٔ دوم
!!a{formula}ی است که همهٔ متغیرهای آن، از جمله متغیرهای محمولی و
تابعی، مقید باشند.
\end{explain}

\begin{defn}[اعتبار]
یک جملهٔ~$!A$ \emph{معتبر} است، یعنی~$\Entails !A$، اگر و تنها اگر
$\Sat{M}{!A}$ برای هر !!{structure}~$\Struct M$ برقرار باشد.
\end{defn}

\begin{defn}[نتیجهٔ منطقی]
مجموعه‌ای از جمله‌ها چون~$\Gamma$ جملهٔ~$!A$ را \emph{نتیجه می‌دهد}،
یعنی~$\Gamma \Entails !A$، اگر و تنها اگر برای هر
!!{structure}~$\Struct M$ که~$\Sat{M}{\Gamma}$، داشته باشیم
$\Sat{M}{!A}$.
\end{defn}

\begin{defn}[ارضاشدنی‌بودن]
مجموعه‌ای از جمله‌ها چون~$\Gamma$ \emph{ارضاشدنی} است هرگاه
$\Sat{M}{\Gamma}$ برای !!{structure}ی چون~$\Struct M$ برقرار باشد. اگر
$\Gamma$ ارضاشدنی نباشد، آن را \emph{ارضاناپذیر} می‌نامیم.
\end{defn}

\end{document}
